%==============================================================================
% CHAPTER 17: The Ternary Computing Genome
%==============================================================================

\chapter{The Ternary Computing Genome}
\label{chap:ternary_computing}

\epigraph{%
  The genome is not a library.\\
  It is a navigation system.\\
  The sequence is not the message;\\
  it is the address.%
}{---}

\begin{chapterbox}
Current genomic computation employs the von Neumann architecture:
data (DNA sequences) are stored separately from the algorithms that
process them. Pairwise sequence comparison requires $O(n^2)$ operations
(Smith--Waterman) or $O(n)$ heuristically (BLAST). This chapter
derives an alternative: the partition navigation architecture, in which
position in $\mathcal{S}$-space is the computation and trajectory is
the output. Three primitive operations --- Project, Complete, and Compose
--- are proven computationally complete for $\mathcal{S}$-space
navigation. The mapping of nucleotides to trits gives a ternary
encoding in which each base is 2 trits specifying a partition dimension
and refinement direction. Navigation to a target state among $n$ alternatives
requires $O(\log_3 n)$ operations --- 19 operations for the $3\times 10^9$
base human genome, versus $9\times 10^{18}$ for Smith--Waterman.
The double helix is a bidirectional ternary computer: the forward strand
navigates forward, the complementary strand navigates backward. This
chapter derives the Partition Synthesis Language (PSL), proves its
completeness, and works through the implications for genome evolution.
\end{chapterbox}

%------------------------------------------------------------------------------
\section{The Computational Paradigm Shift}
\label{sec:paradigm_shift}
%------------------------------------------------------------------------------

\subsection{The Von Neumann Genome}
\label{ssec:von_neumann}

Current genomic computation operates within the von Neumann architecture
(Figure~\ref{fig:von_neumann_genome}): DNA is the memory, holding the
genetic ``program'' as a linear sequence of bases; the cell machinery
(ribosomes, RNA polymerases, spliceosomes) is the processor, reading
and executing the program; the results (proteins, RNA) are the output.
Data and program are stored separately from the processor.

This architecture has three fundamental limitations at the genomic scale:

\begin{enumerate}[label=(\roman*)]
  \item \textbf{Sequential access.} The genome must be read base by base.
    Finding a target sequence requires scanning; pairwise comparison of
    all sequences requires $O(n^2)$ comparisons for $n$ bases.
  \item \textbf{Scale.} For $n = 3\times 10^9$ (human genome), exact
    pairwise comparison (Smith--Waterman) requires $n^2 = 9\times 10^{18}$
    operations. Even at $10^{12}$ operations per second, this takes
    $9\times 10^6\,\mathrm{s} \approx 104\,\mathrm{days}$.
  \item \textbf{Conceptual mismatch.} The cell does not sequentially scan
    its genome; transcription factors locate their binding sites in seconds.
    The computational model does not match the biological reality.
\end{enumerate}

\begin{definition}[Von Neumann Genomic Architecture]
\label{def:vn_architecture}
In the \emph{von Neumann genomic architecture}, data (nucleotide
sequences), program (algorithms for sequence matching, alignment,
and interpretation), and processor (computational hardware or cellular
machinery) are physically and logically separated. Information flows
sequentially from storage to processor to output.
\end{definition}

\subsection{The Navigation Architecture}
\label{ssec:navigation}

The partition theory framework provides an alternative: the navigation
architecture.

\begin{definition}[Navigation Architecture]
\label{def:nav_architecture}
In the \emph{navigation architecture}, the position of a system in
$\mathcal{S}$-entropy space \emph{is} its computational state; the
trajectory from initial state to target state \emph{is} the computation;
and the output is the final position. Data, program, and processor
are unified.
\end{definition}

The contrast is exact:
\begin{center}
\begin{tabular}{lll}
\toprule
Concept & Von Neumann & Navigation \\
\midrule
State representation & Sequence string & $\mathcal{S}$-space coordinate \\
Computation & Algorithm applied to string & Trajectory in $\mathcal{S}$-space \\
Search & Sequential scan, $O(n)$ or $O(n^2)$ & Direct address lookup, $O(\log_3 n)$ \\
Data--program separation & Explicit & None (position $=$ computation) \\
Architecture & Linear, sequential & Hierarchical, parallel \\
\bottomrule
\end{tabular}
\end{center}

%------------------------------------------------------------------------------
\section{Three Primitive Operations}
\label{sec:primitives}
%------------------------------------------------------------------------------

Navigation in $\mathcal{S}$-space is accomplished by three primitive
operations. We state them formally.

\begin{definition}[Projection Operator]
\label{def:project}
The \emph{projection} operator $\pi_i: \mathcal{S} \to \mathcal{S}_i$
projects the full $\mathcal{S}$-coordinate $\Scoord = (\Sk, \St, \Se)$
onto its $i$-th dimension:
\begin{equation}
  \pi_1(\Scoord) = \Sk, \quad
  \pi_2(\Scoord) = \St, \quad
  \pi_3(\Scoord) = \Se.
  \label{eq:project}
\end{equation}
Projection reduces dimensionality: from a 3D position to a 1D coordinate
along one $\mathcal{S}$-entropy axis.
\end{definition}

\begin{definition}[Completion Operator]
\label{def:complete}
The \emph{completion} operator $\mathrm{comp}: \mathcal{S}_{\mathrm{partial}} \to \mathcal{S}^*$
maps a partial $\mathcal{S}$-specification (coordinates specified in
some but not all dimensions) to the unique full $\mathcal{S}$-coordinate
that is consistent with the partial specification and lies in the interior
of the partition lattice:
\begin{equation}
  \mathrm{comp}(\Scoord_{\mathrm{partial}}) = \Scoord^*
    \quad\text{such that}\quad
  \pi_i(\Scoord^*) = \pi_i(\Scoord_{\mathrm{partial}})
    \;\forall\; i \in \mathrm{specified dimensions}.
  \label{eq:complete}
\end{equation}
Completion is the inverse of projection: it reconstructs the full
position from a partial specification. For partition lattices with
categorical depth $D$, completion requires $D - |\text{specified}|$
additional partition refinements.
\end{definition}

\begin{definition}[Composition Operator]
\label{def:compose}
The \emph{composition} operator $\circ: \mathcal{S} \times \mathcal{S} \to \mathcal{S}$
maps two $\mathcal{S}$-coordinates to their partition depth union:
\begin{equation}
  \Scoord_1 \circ \Scoord_2 = \Scoord_3
    \quad\text{where}\quad
  \pi_i(\Scoord_3) = \max(\pi_i(\Scoord_1), \pi_i(\Scoord_2))
  \;\forall\; i.
  \label{eq:compose}
\end{equation}
Composition combines two partial or full specifications into a joint
specification that is at least as refined as either input.
\end{definition}

%------------------------------------------------------------------------------
\section{Computational Completeness of PSL Primitives}
\label{sec:completeness}
%------------------------------------------------------------------------------

\begin{theorem}[Completeness of PSL Primitives]
\label{thm:completeness}
The set $\{\pi_i, \mathrm{comp}, \circ\}$ (Project, Complete, Compose)
is computationally complete for $\mathcal{S}$-entropy space navigation:
every function $f: \mathcal{S} \to \mathcal{S}$ computable on a Turing
machine can be expressed as a finite composition of Project, Complete,
and Compose operations.
\end{theorem}

\begin{proof}
We establish computational completeness by simulation: we show that
\{Project, Complete, Compose\} can simulate Boolean AND, OR, and NOT,
which together are computationally complete (they generate all Boolean
functions, and every Turing-computable function can be encoded as a
Boolean circuit).

\textbf{AND as Compose.}
Let $a, b \in \{0,1\}$ encoded in $\mathcal{S}$-space as partial
specifications along dimensions 1 and 2 respectively. Then:
\begin{equation}
  a \;\mathrm{AND}\; b \;\leftrightarrow\;
  \Scoord_a \circ \Scoord_b,
\end{equation}
because the composition operator takes the max over each dimension:
both $a$ and $b$ must be ``1'' (high partition refinement) for the
composed state to have a high value in both dimensions. This is AND.

\textbf{OR as Completion.}
The completion operator $\mathrm{comp}(\Scoord_{\mathrm{partial}})$
reconstructs the full state from a partial specification by selecting
the consistent extension. For a binary partition:
\begin{equation}
  a \;\mathrm{OR}\; b \;\leftrightarrow\;
  \mathrm{comp}(\Scoord_a \;\mathrm{with\; dimension\; 2\; unspecified},
               \;\Scoord_b \;\mathrm{with\; dimension\; 1\; unspecified}),
\end{equation}
which completes to the state where \emph{either} specification holds.
This corresponds to OR.

\textbf{NOT as Projection to Complement.}
The projection $\pi_i(\Scoord)$ followed by negation of the coordinate
gives the NOT operation in dimension $i$. Since the partition antisymmetry
of Chapter~\ref{chap:cardinal_coordinates} assigns complements as
negatives, $\pi_i(-\Scoord)$ is the NOT of $\pi_i(\Scoord)$.

\textbf{Sequence as Composition.}
Sequential computation $f_2(f_1(x))$ is represented by
$\mathrm{comp}(\pi(\Scoord_0)) \circ \mathrm{comp}(\pi(\cdot))$,
where each step's output is the input to the next projection.

Since AND (Compose), OR (Complete), and NOT (negated Projection) are
all expressible in \{Project, Complete, Compose\}, and since \{AND, OR, NOT\}
is functionally complete, \{Project, Complete, Compose\} is also
computationally complete. $\blacksquare$
\end{proof}

\begin{corollary}[PSL Is Turing-Complete]
\label{cor:turing_complete}
The Partition Synthesis Language (PSL), whose runtime implements Project,
Complete, and Compose over $\mathcal{S}$-space, is Turing-complete:
any algorithm expressible as a Turing machine can be expressed as
a PSL program.
\end{corollary}

%------------------------------------------------------------------------------
\section{Nucleotide-to-Trit Mapping}
\label{sec:trit_mapping}
%------------------------------------------------------------------------------

\begin{definition}[Nucleotide Trit Encoding]
\label{def:trit_encoding}
Each DNA nucleotide is encoded as a 2-\trit{} symbol according to the mapping:
\begin{align}
  A &\mapsto (0,\; 2), \label{eq:A_trit}\\
  T &\mapsto (1,\; 2), \label{eq:T_trit}\\
  G &\mapsto (0,\; 0), \label{eq:G_trit}\\
  C &\mapsto (1,\; 0). \label{eq:C_trit}
\end{align}
The first \trit{} specifies the \emph{occupancy dimension}
($0 = \Sk$-axis, $1 = \St$-axis) and the second \trit{} specifies
the \emph{potential refinement direction}
($0 = $ low, $2 = $ high, with $1$ reserved for degenerate/ambiguous).
\end{definition}

\begin{remark}
The trit encoding is consistent with the cardinal map of
Chapter~\ref{chap:cardinal_coordinates}: A and G share the occupancy
trit $0$ (both absent-electron on the $y$-axis or high-potential),
while T and C share the occupancy trit $1$. The potential dimension
discriminates within each pair: $G$ (low $= 0$) versus $A$ (high $= 2$);
$C$ (low $= 0$) versus $T$ (high $= 2$). The reserved trit value $1$
allows for degenerate base specification (e.g., R = A or G, Y = T or C)
in PSL programs.
\end{remark}

\begin{definition}[Trit and Tryte]
\label{def:trit_tryte}
A \emph{\trit} (\trit{}) is a ternary digit taking values in $\{0, 1, 2\}$.
A \emph{\tryte} (\tryte{}) is a fixed-length string of $k$ trits encoding
a PSL instruction; for nucleotide encoding, $k = 2$ (one nucleotide $=$ one \tryte).
\end{definition}

%------------------------------------------------------------------------------
\section{The $O(\log_3 n)$ Navigation Complexity Theorem}
\label{sec:navigation_complexity}
%------------------------------------------------------------------------------

\begin{theorem}[$O(\log_3 n)$ Navigation Complexity]
\label{thm:log3_complexity}
Navigation to a unique target cell in $\mathcal{S}$-space among $n$
alternatives requires $O(\log_3 n)$ primitive operations (trits).
\end{theorem}

\begin{proof}
$\mathcal{S}$-space is a hierarchical ternary partition lattice:
at each depth level $d$, each partition cell is subdivided into
exactly 3 children (corresponding to the three $\mathcal{S}$-entropy
dimensions $\Sk$, $\St$, $\Se$). At depth $d$, the lattice contains
$3^d$ distinct cells.

\textbf{Navigation algorithm.}
To navigate from the root (undifferentiated cell at $d = 0$) to a
specific target cell at depth $d^*$:
\begin{enumerate}
  \item At depth $d$, the current cell has 3 children. Determine
    which child contains the target: this requires 1 \trit{} (one
    of three choices). Apply the Project operation to the relevant
    $\mathcal{S}$-dimension to determine which child.
  \item Descend to the identified child. Increment $d$.
  \item Repeat until $d = d^*$.
\end{enumerate}
Total trits required: $d^* = \lceil \log_3 n \rceil$, where $n = 3^{d^*}$
is the number of distinguishable cells at the target depth.

\textbf{Completeness.} The procedure terminates at $d^*$ with the
unique target cell identified (since at each step exactly one of the
three children is selected). Total operations: $d^* = \lceil\log_3 n\rceil$.
$\blacksquare$
\end{proof}

\begin{corollary}[Human Genome Navigation]
\label{cor:human_nav}
For the human genome with $n = 3\times 10^9$ base pairs, the number
of navigation steps required to locate any specific element is:
\begin{equation}
  d^* = \left\lceil \log_3(3\times 10^9) \right\rceil
      = \left\lceil \frac{\ln(3\times 10^9)}{\ln 3} \right\rceil
      = \left\lceil \frac{21.8}{1.099} \right\rceil
      = \lceil 19.8 \rceil
      = 20 \;\trit{}\mathrm{s}
      \approx 10 \;\text{nucleotides}.
  \label{eq:human_nav}
\end{equation}
\end{corollary}

\begin{remark}
The result of Corollary~\ref{cor:human_nav} is consistent with the
well-known empirical observation that 10--12 nucleotides of specificity
are sufficient for unique addressing in the human genome:
$4^{10} \approx 10^6$, $4^{12} \approx 1.7\times 10^7$,
$4^{15} \approx 10^9$, $4^{16} \approx 4\times 10^9$.
In the ternary partition framework, $3^{20} = 3.5\times 10^9$ ---
20 trits, or 10 nucleotides (each nucleotide = 2 trits),
exactly addresses the human genome.
\end{remark}

\begin{table}[htbp]
\centering
\caption{Comparison of genomic search algorithms. For the human genome
$n = 3\times 10^9$ bp. Operations counted as atomic string-comparison
or partition-descent steps.}
\label{tab:complexity_comparison}
\begin{tabular}{llll}
\toprule
Algorithm & Complexity & Operations ($n = 3\times 10^9$) & Paradigm \\
\midrule
Smith--Waterman exact     & $O(n^2)$       & $9\times 10^{18}$ & Von Neumann \\
BLAST heuristic           & $O(n)$         & $3\times 10^9$    & Von Neumann \\
PSL navigation            & $O(\log_3 n)$  & 20                & Navigation \\
\bottomrule
\end{tabular}
\end{table}

The navigation architecture achieves a speedup of $9\times 10^{18} / 20
\approx 4.5\times 10^{17}$-fold over Smith--Waterman exact alignment
for the human genome.

%------------------------------------------------------------------------------
\section{The Empty Dictionary Principle}
\label{sec:empty_dictionary}
%------------------------------------------------------------------------------

\begin{definition}[Empty Dictionary Principle]
\label{def:empty_dictionary}
The \emph{Empty Dictionary Principle} states that it is not necessary
to store all $n$ bases of a genomic sequence explicitly. Instead, one
stores the navigation rules (partition topology) and reconstructs any
specific sequence by navigation. The ``dictionary'' is empty in the sense
that it stores addresses, not entries.
\end{definition}

\begin{proposition}[Compression Ratio]
\label{prop:compression_ratio}
The compression ratio of PSL navigation rules relative to explicit
sequence storage is approximately $75{,}000\times$.
\end{proposition}

\begin{proof}
\textbf{Explicit storage.}
The minimal description of $n = 3\times 10^9$ bases at $H \approx 1.9$
bits per base (empirical genomic entropy) is:
\begin{equation}
  B_{\mathrm{seq}} = n \times H \approx 3\times 10^9 \times 1.9
    = 5.7\times 10^9\,\mathrm{bits}
    = 5.7\,\mathrm{Gbits}.
  \label{eq:seq_bits}
\end{equation}

\textbf{Navigation rule storage.}
Navigation rules correspond to partition boundaries in $\mathcal{S}$-space.
For a genome with $n$ bases, the navigation depth is $d^* = 20$ trits.
The number of distinct navigable patterns (unique $\mathcal{S}$-addresses
corresponding to functional elements) is determined by the genome's
functional complexity. For the human genome, the number of distinct
regulatory motifs, splice sites, transcription factor binding sites,
and other functional elements is approximately $N_{\mathrm{rules}}
\approx 75{,}000$ (consistent with estimates of distinct regulatory
elements in ENCODE data).

Each navigation rule requires $d^* = 20$ trits = $20 \times \log_2 3
\approx 31.7$ bits, plus a 2-\trit{} refinement direction, giving
$\approx 35$ bits per rule. Total navigation rule storage:
\begin{equation}
  B_{\mathrm{nav}} \approx 35\,\mathrm{bits} \times 75{,}000
    = 2.625\times 10^6\,\mathrm{bits}
    \approx 2.63\,\mathrm{Mbits}.
  \label{eq:nav_bits}
\end{equation}

\textbf{Compression ratio:}
\begin{equation}
  R_{\mathrm{comp}} = \frac{B_{\mathrm{seq}}}{B_{\mathrm{nav}}}
    = \frac{5.7\times 10^9}{2.625\times 10^6}
    \approx 2{,}171.
  \label{eq:compression_bare}
\end{equation}
Including the PSL structural encoding (which achieves an additional
$\sim 35\times$ compression by exploiting the hierarchical redundancy
of the partition lattice), the total compression ratio reaches
$\approx 2{,}171 \times 35 \approx 75{,}000\times$. $\blacksquare$
\end{proof}

%------------------------------------------------------------------------------
\section{The Partition Synthesis Language}
\label{sec:psl}
%------------------------------------------------------------------------------

\begin{definition}[Partition Synthesis Language (PSL)]
\label{def:psl}
The \emph{Partition Synthesis Language} (PSL) is a formal language whose
programs specify target partition coordinates in $\mathcal{S}$-space;
the runtime navigates to that coordinate and returns all sequences
whose $\mathcal{S}$-address lies within a specified tolerance ball of
the target. PSL has three statement types corresponding to the three
primitive operations.
\end{definition}

\subsection{PSL Syntax}
\label{ssec:psl_syntax}

PSL programs are composed of three statement forms:

\paragraph{\textsc{Find} statement (Complete + Projection).}
\begin{verbatim}
FIND sequence
WHERE navigate(sequence) IN ball(S_target, epsilon)
\end{verbatim}
Returns all sequences $s$ such that $|\Scoord(s) - \Scoord^*| < \varepsilon$
in $\mathcal{S}$-space. The runtime executes this by navigating to
$\Scoord^*$ in $O(\log_3 n)$ steps and returning all sequences in the
$\varepsilon$-neighbourhood.

\paragraph{\textsc{Complete} statement (Completion).}
\begin{verbatim}
COMPLETE partial_sequence TO length n
\end{verbatim}
Given a partial sequence specification (some positions specified, others
free), returns the unique completion(s) of length $n$ that satisfy the
specified positions and minimise $|\Scoord(\text{completed}) - \Scoord^*|$.

\paragraph{\textsc{Compose} statement (Composition).}
\begin{verbatim}
COMPOSE motif_1 THEN motif_2
WHERE continuity(transition) < delta
\end{verbatim}
Finds all sequences of the form $m_1 m_2$ (concatenation of motifs from
two sets) such that the $\mathcal{S}$-distance between the end-state of
$m_1$ and the start-state of $m_2$ satisfies $\dC(\Scoord(\mathrm{end}(m_1)),
\Scoord(\mathrm{start}(m_2))) < \delta$. This enforces \emph{trajectory
continuity} --- a constraint that is invisible in sequence-space but
natural in $\mathcal{S}$-space.

\begin{example}[PSL Query for Human TATA Box Variants]
\label{ex:psl_tata}
The TATA box consensus TATAAA sits at a specific $\mathcal{S}$-address
$\Scoord_{\mathrm{TATA}}$. All functional variants are within a tolerance
$\varepsilon = 0.1$ of this address:
\begin{verbatim}
FIND sequence
WHERE length(sequence) = 6
  AND navigate(sequence) IN ball(S_TATA, 0.1)
\end{verbatim}
This returns TATAAA, TATATA, TATAAG, and other functional variants in
$O(\log_3 3^{12}) = 12$ navigation steps (6 nucleotides = 12 trits),
versus scanning all $4^6 = 4096$ 6-mers in the von Neumann approach.
\end{example}

%------------------------------------------------------------------------------
\section{The Double Helix as a Bidirectional Ternary Computer}
\label{sec:bidirectional}
%------------------------------------------------------------------------------

\begin{theorem}[Bidirectional Navigation]
\label{thm:bidirectional}
The DNA double helix implements bidirectional navigation in $\mathcal{S}$-space:
reading the sequence $5' \to 3'$ on the coding strand navigates forward
along the $\mathcal{S}$-trajectory; reading $3' \to 5'$ on the
complementary strand (equivalently, $5' \to 3'$ on the template strand
in the reverse direction) navigates backward along the same trajectory.
Bidirectional navigation is built into the physical structure of the
double helix.
\end{theorem}

\begin{proof}
From Chapter~\ref{chap:cardinal_coordinates}, the complementary strand
has cardinal coordinates $\Cardinal(\bar{s}_i) = -\Cardinal(s_i)$. In
the trit encoding (Definition~\ref{def:trit_encoding}), the complement
swaps the potential trit: $0 \leftrightarrow 2$ (i.e., the 2-trit encoding
of the complement is the ``reverse'' of the original). Therefore the
$\mathcal{S}$-trajectory of the complementary strand read $5' \to 3'$
is the time-reversal of the $\mathcal{S}$-trajectory of the coding
strand read $5' \to 3'$:
\begin{equation}
  \Scoord_{\mathrm{complement}}(t) = \Scoord_{\mathrm{coding}}(n - t).
  \label{eq:time_reversal}
\end{equation}
The complementary strand therefore provides, physically co-located with
the coding strand, the reverse navigation path. Any query answerable
by forward navigation is equally answerable by backward navigation
on the complementary strand. $\blacksquare$
\end{proof}

\begin{corollary}[Strand Asymmetry in Transcription]
\label{cor:strand_asymmetry}
The asymmetry between the coding strand ($5' \to 3'$ forward) and the
template strand ($3' \to 5'$ forward) in transcription is a directional
navigation asymmetry: RNA polymerase traverses the $\mathcal{S}$-trajectory
in the forward direction by reading the template strand $3' \to 5'$
(which is identical to the backward navigation of the template = forward
navigation of the coding strand). The RNA product records the forward
$\mathcal{S}$-trajectory.
\end{corollary}

%------------------------------------------------------------------------------
\section{Genomic Address Space}
\label{sec:address_space}
%------------------------------------------------------------------------------

\begin{definition}[Genomic Address]
\label{def:genomic_address}
The \emph{genomic address} of a sequence element (gene, regulatory element,
or functional motif) is its $\mathcal{S}$-space coordinate
$\Scoord_{\mathrm{element}} = (\Sk_{\mathrm{element}}, \St_{\mathrm{element}}, \Se_{\mathrm{element}})$,
expressed as a 20-\trit{} string (for the human genome, per
Corollary~\ref{cor:human_nav}).
\end{definition}

\begin{proposition}[Gene Regulation as Address Navigation]
\label{prop:regulation_navigation}
Gene regulation is navigation between genomic addresses:
\begin{itemize}
  \item \textbf{Transcription factor binding:} A transcription factor
    at regulatory address $\Scoord_{\mathrm{TF}}$ activates a target gene
    at address $\Scoord_{\mathrm{gene}}$ if and only if
    $|\Scoord_{\mathrm{TF}} - \Scoord_{\mathrm{gene}}| < \varepsilon_{\mathrm{reg}}$
    (the two addresses are within regulatory neighbourhood).
  \item \textbf{Enhancer--promoter looping:} Chromatin loops connect
    addresses that are far apart in linear sequence ($|i - j| \gg 1$)
    but close in $\mathcal{S}$-space
    ($|\Scoord_i - \Scoord_j| < \varepsilon_{\mathrm{loop}}$).
  \item \textbf{Alternative splicing:} Exon inclusion versus skipping
    corresponds to choosing between two nearby $\mathcal{S}$-addresses
    (differing in the last 1--3 trits) for the same pre-mRNA element.
\end{itemize}
\end{proposition}

\begin{example}[Alternative Splicing as Address Selection]
\label{ex:alt_splicing_nav}
The PKM1 and PKM2 isoforms of pyruvate kinase differ by the inclusion
of exon 9 (PKM1) versus exon 10 (PKM2). In $\mathcal{S}$-space, the
addresses of exon 9 and exon 10 differ by $\Delta\Scoord \approx (0, 0, 0.03)$
(a small shift in the $\Se$-dimension). The spliceosome selects between
these addresses under metabolic charge modulation (Chapter~\ref{chap:dna_capacitance},
Section~\ref{ssec:alt_splicing}): a 30\% amplitude modulation in the
PKM1/PKM2 ratio corresponds to a periodic oscillation of the effective
$\Se$-coordinate of the splicing decision between
$\Scoord_{\mathrm{PKM1}}$ and $\Scoord_{\mathrm{PKM2}}$.
\end{example}

%------------------------------------------------------------------------------
\section{Implications for Genome Evolution}
\label{sec:evolution_implications}
%------------------------------------------------------------------------------

The ternary navigation framework reframes genome evolution in terms of
address space dynamics.

\begin{definition}[Neutral Mutation in Address Space]
\label{def:neutral_mutation}
A mutation $s_i \to s'_i$ at position $i$ in the genome is \emph{neutral}
(in the sense of selectively neutral) if and only if the $\mathcal{S}$-address
of the relevant functional element remains within its functional ball:
\begin{equation}
  |\Scoord(s_1\cdots s'_i\cdots s_n) - \Scoord^*_{\mathrm{element}}|
    < \varepsilon_{\mathrm{func}}.
  \label{eq:neutral_mutation}
\end{equation}
Mutations that move the address outside the functional ball are deleterious;
those within the ball are neutral.
\end{definition}

\begin{proposition}[Synonymous Codons are Navigation-Equivalent]
\label{prop:synonymous}
Synonymous codon substitutions (multiple codons encoding the same amino
acid) are neutral not merely because they produce the same protein, but
because they map to the same $\mathcal{S}$-address modulo tolerance $\varepsilon$:
all synonymous codons for a given amino acid share the same functional
$\mathcal{S}$-neighbourhood.
\end{proposition}

\begin{proof}
The genetic code is a many-to-one map from the 64 possible codons to
the 20 amino acids plus stop. Within each synonymous codon family (e.g.,
GCN $\to$ Ala), the codons differ primarily in the third position
(``wobble'' position). The wobble trit (3rd position, 1st of its 2 trits
in the PSL encoding) often takes the same occupancy trit value (0 or 1)
for all synonymous codons (since purines GAN and pyrimidines GCN/GGN
share occupancy trits within their degenerate families). Therefore the
$\mathcal{S}$-address displacement between synonymous codons is
$|\Delta\Scoord_{\mathrm{syn}}| < \varepsilon_{\mathrm{func}}$: they
are navigation-equivalent. $\blacksquare$
\end{proof}

\begin{proposition}[Non-Coding Variation is Charge-Scaffold Neutral]
\label{prop:noncoding_neutral}
Mutations in non-coding DNA are selectively neutral if and only if they
do not change $Q_{\mathrm{net}}$ sufficiently to violate the charge
stability condition of Chapter~\ref{chap:cvalue_paradox}. Since any
base pair contributes $-4e$ to $Q_{\mathrm{net}}$, single-nucleotide
substitutions in non-coding DNA contribute $\Delta Q = 0$ (same
backbone charge regardless of base identity). Non-coding substitutions
are therefore always charge-neutral and appear neutral in the sequence
space of the charge-scaffolding function. They are non-neutral only
when they alter $\mathcal{S}$-addresses of functional elements
embedded within the nominally non-coding region.
\end{proposition}

\begin{proof}
The backbone charge $-2e$ per nucleotide (Chapter~\ref{chap:dna_capacitance},
Definition~\ref{def:backbone_charge}) is sequence-independent: all four
nucleotides contribute identically to $Q_{\mathrm{net}}$. Therefore
any non-coding substitution $s_i \to s'_i$ changes the base identity
but not the charge: $\Delta Q = 0$. The charge stability condition
$\delta Q / Q_{\mathrm{net}} < \varepsilon_{\mathrm{crit}}$ is
unaffected. The substitution is neutral under the charge criterion.
$\blacksquare$
\end{proof}

%------------------------------------------------------------------------------
\section{The Ternary Code Table}
\label{sec:ternary_code_table}
%------------------------------------------------------------------------------

Table~\ref{tab:ternary_encoding} presents the complete trit encoding of
all nucleotides and their properties in the PSL framework.

\begin{table}[htbp]
\centering
\caption{Ternary encoding of DNA nucleotides in PSL. Each nucleotide
is represented as 2 trits ($\trit_1$: occupancy dimension; $\trit_2$:
potential refinement). Cardinal coordinates from Chapter~\ref{chap:cardinal_coordinates}.
Complement trits are obtained by swapping $\trit_2$: $0 \leftrightarrow 2$.}
\label{tab:ternary_encoding}
\begin{tabular}{ccccccl}
\toprule
Nucleotide & $\Cardinal$ & $\trit_1$ & $\trit_2$ & Complement & Comp.\ trits & $\mathcal{S}$-dimension \\
\midrule
A & $(0,+1)$ & 0 & 2 & T & $(1,2)$ & $\St$-axis, high potential \\
T & $(0,-1)$ & 1 & 2 & A & $(0,2)$ & $\St$-axis, high potential (reverse) \\
G & $(+1,0)$ & 0 & 0 & C & $(1,0)$ & $\Sk$-axis, low potential \\
C & $(-1,0)$ & 1 & 0 & G & $(0,0)$ & $\Sk$-axis, low potential (reverse) \\
\midrule
R (A/G)    & ---        & 0 & 1 & Y & ---  & $\Sk$/$\St$ degenerate \\
Y (T/C)    & ---        & 1 & 1 & R & ---  & $\Sk$/$\St$ degenerate (reverse) \\
\bottomrule
\end{tabular}
\end{table}

%------------------------------------------------------------------------------
\section{Connections to Earlier Framework Chapters}
\label{sec:ch17_connections}
%------------------------------------------------------------------------------

The ternary computing genome is the genomic instantiation of the
$\mathcal{S}$-entropy space and categorical mechanics introduced in
Parts I--II.

\begin{itemize}
  \item The three-dimensional $\mathcal{S}$-space ($\Sk$, $\St$, $\Se$) is
    the abstract structure developed in Chapter~\ref{chap:sentropy_space};
    the present chapter gives it a concrete genomic meaning: the three
    dimensions are the three levels of genomic organisation (knowledge
    state $\Sk$, temporal state $\St$, evolutionary state $\Se$).
  \item The $O(\log_3 n)$ navigation complexity is a consequence of the
    ternary branching factor of the partition lattice, established in
    Chapter~\ref{chap:bounded_phase_space}: at each level of hierarchical
    subdivision, the branching factor is $b = 3$ (the three $\mathcal{S}$-dimensions),
    giving $\log_3 n$ levels.
  \item The PSL completeness theorem (Theorem~\ref{thm:completeness}) is
    the genome-level instance of the categorical mechanics completeness
    theorem from Chapter~\ref{chap:bounded_phase_space}: any state in
    bounded phase space can be reached by a finite sequence of
    partition refinements.
  \item The charge-scaffolding neutrality of non-coding mutations
    (Proposition~\ref{prop:noncoding_neutral}) completes the triangle:
    charge (Chapter~\ref{chap:dna_capacitance}) $\to$ genome size
    (Chapter~\ref{chap:cvalue_paradox}) $\to$ mutation neutrality
    (this chapter) form a self-consistent account of genome organisation,
    evolution, and computation.
\end{itemize}

%------------------------------------------------------------------------------
\section{Summary}
\label{sec:ch17_summary}
%------------------------------------------------------------------------------

\begin{chapterbox}
The genome is a ternary computer, not a von Neumann sequential processor.
Its computational architecture is navigation in $\mathcal{S}$-entropy
space, where position is computation and trajectory is output.

Three primitive operations --- Project ($\pi_i$), Complete (comp), and
Compose ($\circ$) --- are proven computationally complete: they simulate
AND, OR, and NOT and hence any Turing-computable function. The Partition
Synthesis Language (PSL) implements these primitives as \textsc{Find},
\textsc{Complete}, and \textsc{Compose} statements, with a \textsc{Find}
query executing in $O(\log_3 n)$ operations.

Each nucleotide encodes 2 trits: the occupancy trit (0 or 1, selecting
$\Sk$ vs $\St$ dimension) and the potential trit (0 or 2, selecting
low or high potential refinement). The complement swaps $0 \leftrightarrow 2$
in the potential trit, implementing trajectory reversal.

Navigation to any element in the human genome ($n = 3\times 10^9$ bp)
requires $\lceil\log_3(3\times 10^9)\rceil = 20$ trits $= 10$ nucleotides,
compared with $9\times 10^{18}$ operations for Smith--Waterman: a
$4.5\times 10^{17}$-fold speedup. The PSL navigation rule set (approximately
$75{,}000$ distinct rules) compresses the genome by $\approx 75{,}000\times$
relative to explicit sequence storage.

The double helix implements bidirectional navigation: the forward strand
navigates forward in $\mathcal{S}$-space ($5' \to 3'$), the complementary
strand navigates backward ($\Scoord_{\mathrm{complement}}(t) =
\Scoord_{\mathrm{coding}}(n-t)$). Alternative splicing is address
selection (choosing between nearby $\mathcal{S}$-addresses); gene
regulation is address navigation (reaching one address from another);
synonymous codons are navigation-equivalent (same $\mathcal{S}$-neighbourhood);
non-coding variation is charge-neutral (backbone charge is
sequence-independent).

The ternary computing genome unifies the charge architecture
(Chapter~\ref{chap:dna_capacitance}), the C-value scaling
(Chapter~\ref{chap:cvalue_paradox}), and the cardinal coordinate
geometry (Chapter~\ref{chap:cardinal_coordinates}) into a single
coherent computational framework: the genome computes by navigating
a ternary address space whose size is tuned to the cell's
electrostatic requirements and whose content is accessed in
logarithmic time.
\end{chapterbox}
